(1) 一级动力学直接积分给 $N_A(t)=N_0e^{-kt}$、$N_B(t)=N_0(1-e^{-kt})$，转化率 $\xi(t)=1-e^{-kt}$，且 $N_A+N_B=N_0$。这是完全转化且忽略逆反应的理想化模型。

(2) 绝热定容下反应释放的化学能成为显热，$C\dot T=Q(-\dot N_A)=QkN_A$。由初温 $T_0$ 积分得到 $T(t)=T_0+QN_0(1-e^{-kt})/C$，最终 $T_\infty=T_0+QN_0/C$。总化学反应量有限，所以温升有限。

(3) 热释放功率 $P(t)=QkN_0e^{-kt}$，在 $t=0$ 最大值为 $QkN_0$，之后严格下降。累计化学释热 $Q_{\rm rel}(t)=QN_0(1-e^{-kt})=C[T(t)-T_0]$；半温升时 $1-e^{-kt}=1/2$，所以 $t_{1/2}=\ln2/k$。

(4) 若速率依赖当前温度，则 $\dot N_A=-k_0e^{-E_a/(RT)}N_A$，$C\dot T=Qk_0e^{-E_a/(RT)}N_A$。相加 $d[T+(Q/C)N_A]/dt=0$，恒等于 $T_0+QN_0/C$。由于 $k$ 随温度变化，$N_A$ 不再为简单 $e^{-kt}$；需解一维非线性可分离方程或数值积分。
